\chapter{Generating argumentation interpreters as logic programs}
\label{sec:p18-argumentation-meta-interpreters}

Dung's argument-processing unit has two rules: an argument is accepted
when it is not defeated, and it is defeated when an accepted argument
attacks it. The argument-generation unit supplies the attack relation.
This chapter gives the complete program construction and all three
semantic correspondences of the original Section~5
\cite[pp.~348--350, theorem~62]{Dung1995}. The input contract includes
the argument domain, and the well-founded output specifies true,
false and undefined atoms.

\begin{hbdefinition}{def:p18-meta-input}{Argument names, the extensional contract and the source program}
Let \(\mathcal F=(A,R)\) be an arbitrary argumentation framework,
including \(A=\varnothing\). Use the attack direction, image \(S^+\),
defense operator \(\mathsf F\), grounded extension \(G\), and extension
semantics of \cref{def:p18-arbitrary-framework}.
Choose distinct ground names for the arguments and identify \(A\)
with its named image. Let \(D\) be a nonempty Herbrand domain
containing these names. No distinct arguments share a name.

Use unary predicates \(\operatorname{arg},\operatorname{acc},
\operatorname{defeat}\) and one binary predicate
\(\operatorname{atk}\). The argument-generation data consist
of exactly the facts
\[
 \operatorname{arg}(a)\leftarrow\quad(a\in A),\qquad
 \operatorname{atk}(a,b)\leftarrow\quad((a,b)\in R).
\]
There are no other rules defining these predicates.
The relation \(R\) includes all attacks, not merely the attacks
so far discovered by a procedure.

When \(A\ne\varnothing\) and \(D=A\), the source's two-rule
program \(P^0_{\mathcal F}\), with a uniform attack-predicate spelling, is
\[
\begin{aligned}
 \operatorname{acc}(X)&\leftarrow\neg\operatorname{defeat}(X),\\
 \operatorname{defeat}(X)&\leftarrow
                    \operatorname{atk}(Y,X),\operatorname{acc}(Y),
\end{aligned}
\]
together with the attack facts; \(\operatorname{arg}\) is then
unnecessary. Its signature has one constant for every argument,
including isolated arguments, and no function symbols or other
constants. For arbitrary \(A,D\), the guarded program \(P_D\)
uses the extensional facts above and
\[
\begin{aligned}
 \operatorname{acc}(X)&\leftarrow
                \operatorname{arg}(X),\neg\operatorname{defeat}(X),\\
 \operatorname{defeat}(X)&\leftarrow
                \operatorname{arg}(X),\operatorname{arg}(Y),
                \operatorname{atk}(Y,X),\operatorname{acc}(Y).
\end{aligned}
\]
All variables range over \(D\); grounding includes every substitution.
The negative body literal is default negation as in
\cref{def:p18-program-syntax}. The guards make the interpretation
of extra names and the empty framework explicit.
\end{hbdefinition}

\begin{hbdefinition}{def:p18-meta-specialization}{The exact ground specialization and output maps}
Write \(a^\circ=\operatorname{acc}(a)\) and
\(a^\bullet=\operatorname{defeat}(a)\) for \(a\in A\).
Specializing the extensional predicates gives the ground program
\[
 Q_{\mathcal F}=
 \{a^\circ\leftarrow\neg a^\bullet:a\in A\}
 \ \cup\
 \{b^\bullet\leftarrow a^\circ:(a,b)\in R\}
\]
on the specified base \(H_{\mathcal F}=\{a^\circ,a^\bullet:a\in A\}\).
When \(A=\varnothing\), this is the empty program on the empty base.
It is a ground semantic presentation; an arbitrary infinite set
of arguments is not asserted to have a finite source encoding.

For \(S\subseteq A\), set
\[
\begin{aligned}
 m_Q(S)&=\{a^\circ:a\in S\}\cup\{b^\bullet:b\in S^+\},\\
 L_Q(S)&=m_Q(S)\cup\{\neg a^\bullet:a\in S\}
                         \cup\{\neg b^\circ:b\in S^+\}.
\end{aligned}
\]
Let \(B_D^+\) be the set of extensional facts, viewed as atoms.
Define the complete fixed background
\[
\begin{aligned}
 B_D={}&B_D^+
 \cup\{\neg\operatorname{arg}(d):d\in D\setminus A\}\\
 &{}\cup\{\neg\operatorname{atk}(c,d):(c,d)\in D^2\setminus R\}\\
 &{}\cup\{\neg\operatorname{acc}(d),\neg\operatorname{defeat}(d):
                                                  d\in D\setminus A\}.
\end{aligned}
\]
For a literal set \(I\) on \(H_{\mathcal F}\), its lift is
\(\iota_D(I)=B_D\cup I\). Put
\[
             m_D(S)=B_D^+\cup m_Q(S),\qquad
             L_D(S)=\iota_D(L_Q(S)).
\]
The signed sets \(L_Q(S),L_D(S)\) are partial interpretations
exactly when \(S\cap S^+=\varnothing\).
For the source program on \(D=A\), omit the
\(\operatorname{arg}\) predicate and its literals from these formulas.
Denote the resulting positive and signed maps by \(m_0(S)\)
and \(L_0(S)\), respectively: they are obtained from \(m_A(S)\)
and \(L_A(S)\) by deleting all domain literals.
\end{hbdefinition}

All reducts, positive least models, unfounded sets, and Fitting
operators below use the exact definitions and arbitrary-ground-program
domain of \cref{def:p18-program-reduct,def:p18-unfounded-operators,def:p18-completion-semantics}. In particular
\(\Gamma_P(X)=\operatorname{LM}(P^X)\),
\(\Phi_P(I)=T_P(I)\cup\neg Z_P(I)\), and
\(V_P(I)=T_P(I)\cup\neg U_P(I)\).

\begin{hblemma}{thm:p18-meta-reduct}{proved=annex:p18.meta.reduct}{The complete reduct calculation}
For any set \(X\) of ground atoms of the appropriate program, let
\[
                         d_X=\{a\in A:a^\bullet\in X\}.
\]
Then
\[
\begin{aligned}
 \Gamma_Q(X)&=m_Q(A\setminus d_X),\\
 \Gamma_{P_D}(X)&=m_D(A\setminus d_X).
\end{aligned}
\]
For \(A\ne\varnothing\), also
\(\Gamma_{P^0_{\mathcal F}}(X)=m_0(A\setminus d_X)\).
These formulas cover every atom, not just the accepted arguments.
If \(N(S)=A\setminus S^+\), then, for each corresponding output map \(m\),
\[
                   \Gamma(m(S))=m(N(S)),\qquad
                   \Gamma^2(m(S))=m(\mathsf F(S)).
\]
All equalities allow infinitely many attackers.
\end{hblemma}

\begin{hbtheorem}{thm:p18-meta-stable}{proved=annex:p18.meta.stable}{Both directions of the stable-model correspondence}
The maps
\[
                   E\longmapsto m_Q(E),\qquad
                   E\longmapsto m_D(E)
\]
are bijections from stable extensions of \(\mathcal F\) onto
stable models of \(Q_{\mathcal F}\) and \(P_D\), respectively.
For \(A\ne\varnothing\), \(E\mapsto m_0(E)\) is the same
bijection for the literal source program \(P^0_{\mathcal F}\).
The inverse extracts the arguments whose acceptance atoms are true.
This includes an empty extension, self-attacks and infinite frameworks.
For a stable \(E\), \(L_Q(E)\) and \(L_D(E)\) are exactly the full
total interpretations determined by those positive stable models.
The analogous total interpretation for the source program is \(L_0(E)\).
\end{hbtheorem}

\begin{hblemma}{thm:p18-meta-specialization}{proved=annex:p18.meta.specialization}{Specialization preserves the full partial semantics}
A literal set \(J\) is a consistent fixed point of \(\Phi_{P_D}\)
if and only if
\[
                     J=\iota_D(I),\qquad I=\Phi_Q(I)
\]
for a consistent \(I\) on \(H_{\mathcal F}\).
The maps are inverse order isomorphisms. Consequently
\[
                       \operatorname{FM}(P_D)
                                      =\iota_D(\operatorname{FM}(Q)).
\]
The same assertion holds for the source program on \(D=A\), with
the stated omission of domain literals.
Thus replacing an extensional attack body by its actual edge
instances is a proved specialization of these programs.
\end{hblemma}

\begin{hblemma}{thm:p18-meta-defense-steps}{proved=annex:p18.meta.defense.steps}{Four precise Fitting steps for one defense step}
For every \(S\subseteq A\), on the auxiliary literal powerset,
\[
                             \Phi_Q^4(L_Q(S))=L_Q(\mathsf F(S)).
\]
Let \(S_0=\varnothing\), \(S_{\alpha+1}=\mathsf F(S_\alpha)\),
with unions at limits. Define
\[
 \tau_0=0,\qquad \tau_{\alpha+1}=\tau_\alpha+4,\qquad
 \tau_\lambda=\sup_{\beta<\lambda}\tau_\beta
                   \quad(\lambda\ne0\text{ a limit}).
\]
Then the actual Fitting iteration satisfies
\[
                        \Phi_Q^{\tau_\alpha}(\varnothing)=L_Q(S_\alpha).
\]
The iteration on the left uses unions at limit ordinals.
In particular \(\operatorname{FM}(Q)=L_Q(G)\).
\end{hblemma}

\begin{hbtheorem}{thm:p18-meta-completion}{proved=annex:p18.meta.completion}{Complete extensions are exactly the partial completion models}
The map \(E\mapsto L_Q(E)\) is an order isomorphism between
the complete extensions of \(\mathcal F\) and the consistent
fixed points of \(\Phi_Q\), with inverse the true acceptance set.
Together with \(\iota_D\), it gives the same correspondence
for \(P_D\). It sends grounded extensions to the least partial
completion models, preferred extensions to maximal partial
completion models, and stable extensions to total completion
models, in both directions.

For infinite input this refers to the head-by-head operator
semantics of completion in \cref{def:p18-completion-semantics};
it does not replace an infinite disjunction by a finitary formula.
The truth-value equality convention for partial completion is
retained.
\end{hbtheorem}

\begin{hblemma}{thm:p18-meta-positive-rank}{proved=annex:p18.meta.positive.rank}{Why well-founded and Fitting semantics agree here}
Suppose a ground normal program has a rank on its atoms in
\(\{0,\ldots,r\}\), for a nonnegative integer \(r\), such that every positive body atom has
strictly smaller rank than its head. Its operators \(\Phi\)
and \(V\) have exactly the same fixed points on the auxiliary
literal powerset, and hence the same least fixed point.

This applies to \(Q_{\mathcal F}\), with acceptance rank 0
and defeat rank 1. It applies to \(P_D\), with domain and
attack rank 0, acceptance rank 1 and defeat rank 2.
It also applies to the source program, with attack and
acceptance rank 0 and defeat rank 1. Therefore
\[
                           \operatorname{WFM}(P)=\operatorname{FM}(P)
\]
for each of these three encodings, even if \(A,D,R\) are infinite.
No acyclicity of the original attack graph is assumed.
\end{hblemma}

\begin{hbtheorem}{thm:p18-meta-full-correspondence}{proved=annex:p18.meta.full.correspondence}{The complete signed form of Dung's interpreter theorem}
For every \(E\subseteq A\),
\[
\begin{aligned}
 E\text{ is grounded}
 &\quad\Longleftrightarrow\quad L_Q(E)=\operatorname{WFM}(Q)\\
 &\quad\Longleftrightarrow\quad L_D(E)=\operatorname{WFM}(P_D).
\end{aligned}
\]
For \(A\ne\varnothing\), these are also equivalent to
\(L_0(E)=\operatorname{WFM}(P^0_{\mathcal F})\).
In these well-founded models the accepted arguments are exactly \(G\);
the defeated arguments are exactly \(G^+\). Their dynamic literals are
\[
       a^\circ,\neg a^\bullet\quad(a\in G),\qquad
       b^\bullet,\neg b^\circ\quad(b\in G^+).
\]
For \(c\in A\setminus(G\cup G^+)\), both \(c^\circ\) and
\(c^\bullet\) are undefined. The guarded model's background
literals are exactly \(B_D\); the source program omits domain
literals, and the reduced program has no background predicates.
Parts (1)--(3) of the source theorem are thus supplied by the
stable-model bijection, this full grounded equivalence, and
the well-founded/Fitting equality.

If the programs are themselves turned into Dung's literal
support or clause argument graphs, then
\[
\begin{aligned}
 \operatorname{con}(G_{\mathcal A_\infty(P_D)})
  &=L_D(G)
   =\operatorname{con}(G_{\mathcal A_{\rm c}(P_D)}).
\end{aligned}
\]
The support graph retains every finite support and every negative
unit argument from \cref{def:p18-source-program-graphs}.
It is not replaced by the original framework \(\mathcal F\).
Writing \(H_{P_D}\) for the full ground atom base, its stable
support extensions are exactly
\[
       {\cal E}_\infty(H_{P_D}\setminus m_D(E))
                         \qquad(E\in\operatorname{Stab}(\mathcal F)).
\]
The inverse extracts true acceptance atoms from the positive conclusions.
\end{hbtheorem}

\begin{hbmethod}{def:p18-meta-finite-procedure}{Constructing and evaluating a finite interpreter}
Input a complete finite list of \(n\) argument names and \(m\)
directed attacks. Emit one acceptance clause per argument and
one defeat clause per edge, obtaining \(Q_{\mathcal F}\) with
\(n+m\) clauses and \(2n\) atoms. Keep isolated arguments in the list.
Index the names by \(1,\ldots,n\); the bounds below count operations
on these indices and assume constant-time membership in finite Boolean arrays.

To compute the grounded answer, iterate \(\mathsf F\) from the empty
set and then output \(L_Q(G)\), or its full lift \(L_D(G)\).
For each step, scan the edges once to form \(S^+\); initialize
every argument as defended, and reject a target \(b\) whenever
an edge \(a\to b\) has \(a\notin S^+\).
To enumerate stable answers, enumerate all \(E\subseteq A\)
and retain precisely those with \(E\cap E^+=\varnothing\)
and \(E\cup E^+=A\); output their \(m_Q(E)\) or \(m_D(E)\).
\end{hbmethod}

\begin{hbtheorem}{thm:p18-meta-finite-procedure}{proved=annex:p18.meta.finite.procedure}{Termination and the exact finite input bounds}
The stated construction and procedures are correct.
The grounded procedure uses at most \(n+1\) defense evaluations
and time \(O((n+1)(n+m+1))\) with explicit finite adjacency data.
The stable enumeration uses time \(O(2^n(n+m+1))\).
Construction of the sparse dynamic program uses \(O(n+m)\)
output space. These bounds concern the supplied finite framework
and dynamic output; if all negative attack literals are also
listed for \(D=A\), the background output can require \(O(n^2)\)
space.

Infinite ground programs retain the semantic theorems. They do
not inherit these terminating finite enumeration procedures.
Effective generation of arguments and attacks is a separate
input obligation of an argument-generation unit.
\end{hbtheorem}

\begin{hbtheorem}{thm:p18-meta-boundary-examples}{proved=annex:p18.meta.boundary.examples}{Exact examples of domain and semantic failures}
The following consequences follow from the declared encodings.
\begin{enumerate}
\item For \(A=\{a,b\}\), \(R=\{(a,b)\}\), the grounded and sole
      stable extension is \(\{a\}\). The dynamic well-founded
      model is
      \[
              \{a^\circ,\neg a^\bullet,b^\bullet,\neg b^\circ\}.
      \]
      A purported full model omitting \(\neg b^\circ\) is incomplete.
\item For an isolated \(a\), the unique dynamic answer is
      \(\{a^\circ,\neg a^\bullet\}\). Inferring argument names
      solely from attack facts loses this argument.
\item On a domain with an extra name \(d\notin A\), omitting the
      domain guard from the acceptance rule makes acceptance of \(d\)
      true in the well-founded and Fitting models.
      For \(A=\varnothing,D=\{d\}\), this
      happens even though the framework has no argument at all.
      The guarded program makes both acceptance and defeat of
      \(d\) false.
\item For a single self-attacking argument, there is no stable
      model, while the dynamic well-founded and Fitting models
      leave both atoms undefined. For a two-argument mutual attack
      there are two stable models, corresponding to \(\{a\}\)
      and \(\{b\}\), while the dynamic grounded answer is empty.
\item Using ordinary classical models of the program implications
      gives a different result: for an isolated \(a\), the
      assignment \(a^\circ=\mathrm{false}\),
      \(a^\bullet=\mathrm{true}\) satisfies those implications
      but is not a stable or completion model.
\item If the generator writes \(\operatorname{attacks}(a,b)\)
      while the processing rule reads a distinct predicate
      \(\operatorname{atk}\), the one-edge input in (1) produces
      acceptance of both arguments. A uniform predicate signature
      is essential.
\end{enumerate}
In the literal support graph even an acceptance conclusion can
attack an argument: a support proof for \(a^\circ\) attacks
the unit \((\{\neg a^\circ\},\neg a^\circ)\).
Thus the program-to-argumentation bridge must include all units
and their attacks.
\end{hbtheorem}

\begin{hbtheorem}{thm:p18-meta-infinite-input}{proved=annex:p18.meta.infinite.input}{Decidable infinite attack data need not give decidable answers}
There is a countable framework with decidable argument names
and a decidable attack relation, of incoming-path length at
most one, for which grounded acceptance is undecidable.
Its grounded extension is its unique stable extension, and
its generated well-founded model is total on the dynamic atoms.
Thus the difficulty persists even with agreement of all these
semantics and with no cycles.
\end{hbtheorem}

The generator/processor architecture specifies a semantics for the
framework supplied to it. It does not establish the truth of the
arguments' premises or complete an unfinished enumeration of attacks.
The edge specialization above is one fully proved program
transformation. The source's further suggestions about partial
evaluation are implementation directions; they do not assert that
every transformation, search order, or interpreter terminates.
